\documentclass[10pt,a4paper]{amsart}
\usepackage[margin=19mm]{geometry}
\usepackage{amsmath,amssymb,mathtools}
\usepackage[scr=rsfs,cal=euler]{mathalfa}
\usepackage{xfrac}
\usepackage[unicode,hidelinks,bookmarks=false]{hyperref}
\newcommand{\ie}{i.e.\ }
\newcommand{\cf}{cf.\ }
\newcommand{\resp}{respectively\ }
\newcommand{\Subsection}{\subsection}
\providecommand{\nobreakdash}{\nobreak\hbox{-}}
\setlength{\emergencystretch}{2em}
\multlinegap0pt
\usepackage{bm}
\usepackage{bbm}
\usepackage{dsfont}
\usepackage{xspace}
\usepackage{cancel}
\usepackage{mathtools}
\usepackage{amsthm}
\makeatletter
\@ifundefined{theorem}{\newtheorem{theorem}{Theorem}}{}
\@ifundefined{remark}{\newtheorem{remark}[theorem]{Remark}}{}
\makeatother
\title[The aspect ratio of the Twin Dragon is $1/\phi$]{The aspect ratio of the Twin Dragon is $1/\phi$}
\author{Dmitry Mekhontsev}
\date{Source: arXiv:2604.05010v5 (CC BY 4.0)}
\begin{document}
\maketitle

\noindent\emph{Source and licence.} The aspect ratio of the Twin Dragon is $1/\phi$. Version arXiv:2604.05010v5 (CC BY 4.0), distributed under the Creative Commons Attribution 4.0 International licence, as recorded in the arXiv licence metadata for this exact version and shown on the abstract page https://arxiv.org/abs/2604.05010v5. Full rights notice: licenses/arxiv-2604.05010-CC-BY-4.0.txt.\par\smallskip

\noindent\textit{Editorial scope.} Counted unit: the Theorem whose source label is \texttt{thm:family} in the article (source0.tex lines 195--227, source label \texttt{thm:family}), delivered together with the article's own complete original proof of it (source0.tex lines 229--280, one \texttt{proof} environment of 52 lines). No mathematics is written, repaired or re-derived: statement and proof are the article's own text line for line. Required context retained uncounted: eq:psi-simple (lines 301--303); rem:orientation (lines 605--619). Every definition, display and label of those lines is retained unchanged. Omitted from this package and not redistributed: 1--194, 228--228, 281--300, 304--604, 620--1209. Nothing inside a retained excerpt is omitted or altered: every retained body is delivered exactly as the article's lines are written, so no omission receipt of any kind accompanies this package. Citations: the retained excerpts carry no citation command at all, so no bibliography entry of the article is transcribed and no reference is redistributed. Reference adaptation: none was needed and none was made; every reference inside the delivered unit resolves inside it, so no unresolved reference or citation is produced. Printed slips kept literal: no printed word, symbol, hypothesis or number of the counted statement or of its proof was altered. Layout and class: the article's own class is \texttt{amsart} with packages including amsmath, amssymb, amsthm, enumerate, graphicx, hyperref, pgfplots; the wrapper is the lane's pinned \texttt{amsart} editorial wrapper at 10pt on A4, which restates the theorem-like environments the retained text needs and adds the article's own macro definitions that the retained text needs (none), with extra wrapper packages (bm, bbm, dsfont, xspace, cancel, mathtools). The delivered numbering is the unit's own. Assets: no retained span contains a figure environment, a graphics-inclusion command, a PDF-page inclusion, a table or a TikZ picture; no image file is copied into this package, the package declares no asset dependency and no image file is redistributed with it. Licence: version arXiv:2604.05010v5 of this article is distributed under the Creative Commons Attribution 4.0 International licence, as recorded in the arXiv licence metadata for this exact version and shown on the abstract page https://arxiv.org/abs/2604.05010v5. A full rights notice is delivered at licenses/arxiv-2604.05010-CC-BY-4.0.txt. Characters: every retained excerpt is pure ASCII, so the delivered engine, XeLaTeX with its default Latin Modern text font, reports no missing character.
\begin{equation}\label{eq:psi-simple}
  \tan 2\psi = \frac{\operatorname{Im}(a^2)}{\operatorname{Re}(a^2)-|a|^4}\,;
\end{equation}

\begin{remark}\label{rem:orientation}
The orientation formula~\eqref{eq:psi} follows from the decomposition
of the covariance matrix.  Writing $\omega = |\omega|\, e^{2i\psi}$
for $Z = X + iY$ with $E[Z]=0$:
\[
  M = \frac{E[|Z|^2]}{2}\,I
    + \frac{|\omega|}{2}
      \begin{pmatrix}\cos 2\psi & \sin 2\psi \\ \sin 2\psi & -\cos 2\psi\end{pmatrix},
\]
whose larger eigenvalue has eigenvector $(\cos\psi, \sin\psi)$.
For the collinear case~\eqref{eq:psi-simple},
$\operatorname{Im}(a^2)$ measures the rotation per level of the
IFS hierarchy, and $\operatorname{Re}(a^2)-|a|^4$ is the real drift;
their ratio determines the principal-axis direction.
\end{remark}

\begin{theorem}\label{thm:family}
With the notation above, let $I$ and~$J$ denote the sets of
orientation-preserving and -reversing indices, and define
\[
  u = \frac{1}{N}\!\sum_{k \in I}s_k^2,\qquad
  v = \frac{1}{N}\!\sum_{k \in J}s_k^2.
\]
Then the second moment $\omega := E_\mu[Z^2]$ is
\begin{equation}\label{eq:omega}
  \omega
  = \frac{a^2\bigl(\tau + \overline{a^2}(v\,\overline{\tau}
          - \bar{u}\,\tau)\bigr)}
         {|1-a^2 u|^2 - |a|^4|v|^2},
\end{equation}
the aspect ratio is
$\mathrm{AR}^2 = (1-\rho)/(1+\rho)$, where
\begin{equation}\label{eq:rho}
  \rho = \frac{(1-|a|^2)\,|\omega|}{|a|^2\,\sigma^2},
\end{equation}
and the major axis of the attractor makes angle
\begin{equation}\label{eq:psi}
  \psi = \tfrac{1}{2}\arg\omega
\end{equation}
with the positive $x$-axis.
When $v = 0$ \textup{(}all maps orientation-preserving\textup{)},
\eqref{eq:omega} reduces to
$\omega = a^2\tau/(1-a^2 u)$, and
\begin{equation}\label{eq:rho-pres}
  \rho = \frac{(1-|a|^2)\,|\tau|}{|1-a^2 u|\,\sigma^2}.
\end{equation}
In general, the aspect ratio depends on $u$, $v$, and the
complex value of~$\tau$ \textup{(}not just~$|\tau|$\textup{)}.
\end{theorem}

\begin{proof}
Let $Z$ be distributed according to the invariant measure $\mu$.
The fixed-point equation
\begin{equation}\label{eq:fp}
  E[h(Z)] = \frac{1}{N}\!\Bigl(
    \sum_{k\in I} E\bigl[h\bigl(a(s_k Z + t_k)\bigr)\bigr]
    + \sum_{k\in J} E\bigl[h\bigl(a(s_k\overline{Z} + t_k)\bigr)\bigr]
  \Bigr)
\end{equation}
applied to $h(z)=z$ gives
$E[Z] = a\bigl(u'\,E[Z] + v'\,\overline{E[Z]}
  + \tfrac{1}{N}\!\sum t_k\bigr)$
where $u' = \frac{1}{N}\sum_{k\in I}s_k$,
$v' = \frac{1}{N}\sum_{k\in J}s_k$;
since $|a|(|u'|+|v'|)\leq|a|<1$ and $\sum t_k = 0$,
the unique solution is $E[Z]=0$.

Applying \eqref{eq:fp} to $h(z)=|z|^2$:
since $E[Z]=0$, the cross terms $2\operatorname{Re}(\overline{s_k Z}\,t_k)$
vanish in expectation, and $|s_k z|^2 = |z|^2$, giving
\[
  E[|Z|^2] = |a|^2\bigl(E[|Z|^2] + \sigma^2\bigr),
\]
so $E[|Z|^2] = |a|^2\sigma^2/(1-|a|^2)$.

Applying \eqref{eq:fp} to $h(z)=z^2$ and writing $\omega := E[Z^2]$:
$(s_k z)^2 = s_k^2 z^2$ for preserving maps and
$s_k^2\overline{z}^2$ for reversing ones, while the cross
terms $2s_k z\cdot t_k$ vanish in expectation
and the translations contribute $\sum t_k^2 = N\tau$, so
\begin{equation}\label{eq:omega-fp}
  \omega = a^2\bigl(u\,\omega + v\,\overline{\omega} + \tau\bigr).
\end{equation}
Rewriting as the $\mathbb{R}$-linear system
$(1-a^2 u)\,\omega - a^2 v\,\overline{\omega} = a^2\tau$
and applying Cramer's rule
($\Delta := |1-a^2 u|^2 - |a|^4|v|^2 > 0$
since $|a|^2(|u|+|v|)\leq |a|^2 < 1$)
gives~\eqref{eq:omega}.
The ellipticity
$\rho = |\omega|/E[|Z|^2]$ gives~\eqref{eq:rho};
since the covariance matrix $M$ of the real 2D vector
$(\mathrm{Re}\,Z,\mathrm{Im}\,Z)$ is positive semidefinite,
$|\omega|\leq E[|Z|^2]$ and $\rho\in[0,1]$, so
$\mathrm{AR}^2 = (1-\rho)/(1+\rho)\in[0,1]$.
The orientation~\eqref{eq:psi} follows from the
covariance decomposition in Remark~\ref{rem:orientation}.

When $v=0$, the equation is $\mathbb{C}$-linear:
$\omega = a^2\tau/(1-a^2 u)$,
giving~\eqref{eq:rho-pres}.
\end{proof}

\end{document}
